\documentclass[10pt,a4paper]{amsart}
\usepackage{amsmath,amssymb,amsthm,graphicx}
\usepackage{verbatim}
\newenvironment{reference}{\comment}{\endcomment}
%\newenvironment{reference}{}{}
\newenvironment{slogan}{\comment}{\endcomment}
\newenvironment{history}{\comment}{\endcomment}

% For commutative diagrams we use Xy-pic
\usepackage[all]{xy}

% We use 2cell for 2-commutative diagrams.
\xyoption{2cell}
% Theorem environments.
%
\theoremstyle{plain}
\newtheorem{theorem}[subsection]{Theorem}
\newtheorem{proposition}[subsection]{Proposition}
\newtheorem{lemma}[subsection]{Lemma}

\theoremstyle{definition}
\newtheorem{definition}[subsection]{Definition}
\newtheorem{example}[subsection]{Example}
\newtheorem{exercise}[subsection]{Exercise}
\newtheorem{situation}[subsection]{Situation}

\theoremstyle{remark}
\newtheorem{remark}[subsection]{Remark}
\newtheorem{remarks}[subsection]{Remarks}

\numberwithin{equation}{subsection}

% Macros
%
\def\lim{\mathop{\mathrm{lim}}\nolimits}
\def\colim{\mathop{\mathrm{colim}}\nolimits}
\def\Spec{\mathop{\mathrm{Spec}}}
\def\Hom{\mathop{\mathrm{Hom}}\nolimits}
\def\Ext{\mathop{\mathrm{Ext}}\nolimits}
\def\SheafHom{\mathop{\mathcal{H}\!\mathit{om}}\nolimits}
\def\SheafExt{\mathop{\mathcal{E}\!\mathit{xt}}\nolimits}
\def\Sch{\mathit{Sch}}
\def\Mor{\mathop{\mathrm{Mor}}\nolimits}
\def\Ob{\mathop{\mathrm{Ob}}\nolimits}
\def\Sh{\mathop{\mathit{Sh}}\nolimits}
\def\NL{\mathop{N\!L}\nolimits}
\def\CH{\mathop{\mathrm{CH}}\nolimits}
\def\proetale{{pro\text{-}\acute{e}tale}}
\def\etale{{\acute{e}tale}}
\def\QCoh{\mathit{QCoh}}
\def\Ker{\mathop{\mathrm{Ker}}}
\def\Im{\mathop{\mathrm{Im}}}
\def\Coker{\mathop{\mathrm{Coker}}}
\def\Coim{\mathop{\mathrm{Coim}}}

% Boxtimes
%
\DeclareMathSymbol{\boxtimes}{\mathbin}{AMSa}{"02}

%
% Macros for moduli stacks/spaces
%
\def\QCohstack{\mathcal{QC}\!\mathit{oh}}
\def\Cohstack{\mathcal{C}\!\mathit{oh}}
\def\Spacesstack{\mathcal{S}\!\mathit{paces}}
\def\Quotfunctor{\mathrm{Quot}}
\def\Hilbfunctor{\mathrm{Hilb}}
\def\Curvesstack{\mathcal{C}\!\mathit{urves}}
\def\Polarizedstack{\mathcal{P}\!\mathit{olarized}}
\def\Complexesstack{\mathcal{C}\!\mathit{omplexes}}
% \Pic is the operator that assigns to X its picard group, usage \Pic(X)
% \Picardstack_{X/B} denotes the Picard stack of X over B
% \Picardfunctor_{X/B} denotes the Picard functor of X over B
\def\Pic{\mathop{\mathrm{Pic}}\nolimits}
\def\Picardstack{\mathcal{P}\!\mathit{ic}}
\def\Picardfunctor{\mathrm{Pic}}
\def\Deformationcategory{\mathcal{D}\!\mathit{ef}}
\usepackage{cleveref}
\crefname{originalequation}{original equation}{original equations}
\Crefname{originalequation}{Original equation}{Original equations}
\makeatletter
\newlabel{equation-epsilon}{{08H7}{source}{}{sourcepointer-1}{}}
\newlabel{equation-epsilon@cref}{{[originalequation][0][]08H7}{[1][0][]source}{}{}{}}
\makeatother
\crefname{originalreference}{original reference}{original references}
\Crefname{originalreference}{Original reference}{Original references}
\makeatletter
\newlabel{lemma-descend-finite-type}{{08HD}{source}{}{sourcepointer-2}{}}
\newlabel{lemma-descend-finite-type@cref}{{[originalreference][0][]08HD}{[1][0][]source}{}{}{}}
\makeatother
\crefname{originalreference}{original reference}{original references}
\Crefname{originalreference}{Original reference}{Original references}
\makeatletter
\newlabel{lemma-epsilon-flat}{{08H8}{source}{}{sourcepointer-3}{}}
\newlabel{lemma-epsilon-flat@cref}{{[originalreference][0][]08H8}{[1][0][]source}{}{}{}}
\makeatother
\crefname{originaldefinition}{original definition}{original definitions}
\Crefname{originaldefinition}{Original definition}{Original definitions}
\makeatletter
\newlabel{sites-cohomology-definition-pseudo-coherent}{{08FT}{source}{}{sourcepointer-4}{}}
\newlabel{sites-cohomology-definition-pseudo-coherent@cref}{{[originaldefinition][0][]08FT}{[1][0][]source}{}{}{}}
\makeatother
\crefname{originallemma}{original lemma}{original lemmas}
\Crefname{originallemma}{Original lemma}{Original lemmas}
\makeatletter
\newlabel{sites-cohomology-lemma-cone-pseudo-coherent}{{08FV}{source}{}{sourcepointer-5}{}}
\newlabel{sites-cohomology-lemma-cone-pseudo-coherent@cref}{{[originallemma][0][]08FV}{[1][0][]source}{}{}{}}
\makeatother
\crefname{originallemma}{original lemma}{original lemmas}
\Crefname{originallemma}{Original lemma}{Original lemmas}
\makeatletter
\newlabel{sites-cohomology-lemma-finite-cohomology}{{08FX}{source}{}{sourcepointer-6}{}}
\newlabel{sites-cohomology-lemma-finite-cohomology@cref}{{[originallemma][0][]08FX}{[1][0][]source}{}{}{}}
\makeatother
\crefname{originallemma}{original lemma}{original lemmas}
\Crefname{originallemma}{Original lemma}{Original lemmas}
\makeatletter
\newlabel{sites-cohomology-lemma-pseudo-coherent-pullback}{{08H4}{source}{}{sourcepointer-7}{}}
\newlabel{sites-cohomology-lemma-pseudo-coherent-pullback@cref}{{[originallemma][0][]08H4}{[1][0][]source}{}{}{}}
\makeatother
\title{Pseudo-coherence descends from the small etale site}
\author{The Stacks Project Authors}
\date{stacks-project:spaces-perfect.tex:local-snapshot-20261012; GNU FDL 1.2 or later}
\begin{document}
\maketitle
\noindent\textit{Editorial note.} All delivered mathematical text is the author's own, copied from the pinned Stacks Project chapter file spaces-perfect.tex; no statement or proof step was rewritten, reordered or removed.\par\medskip
\noindent\textit{Editorial note.} The counted proof is the author's complete argument: load-bearing boundedness input is the delivered lemma-epsilon-flat (L48088)\par\medskip
\noindent\textit{Editorial note.} The label lemma-descend-pseudo-coherent names the counted statement, and the proof environment immediately following it in the source is delivered verbatim as its proof.\par\medskip
\noindent\textit{Editorial note.} The counted statement is carried by the author's own lemma environment in the chapter source.\par\medskip
\noindent\textit{Allowed standard prerequisites (editorial).}\begin{itemize}
\item Lemma lemma-epsilon-flat (delivered, L48088): the comparison functor epsilon\textasciicircum{}* is flat, identifying boundedness of E with boundedness of epsilon\textasciicircum{}*E
\item Lemma lemma-descend-finite-type (in-source): finiteness of cohomology modules descends from the small etale site
\item Cohomology on Sites, Lemmas sites-cohomology-lemma-pseudo-coherent-pullback, sites-cohomology-lemma-finite-cohomology, sites-cohomology-lemma-cone-pseudo-coherent
\end{itemize}\medskip
\ifdefined\hypertarget\hypertarget{sourcepointer-1}{}\fi
\noindent\textit{Source reference.} Original equation 08H7: sources/stacks-spaces-perfect/pinned-original-source.tex, source lines 320--320. Prerequisite pointer invoked at 3140; it names a standard fact of the same chapter that the counted argument uses, and every step of the counted proof built on it is retained verbatim in the delivered body.\par\medskip
\ifdefined\hypertarget\hypertarget{sourcepointer-2}{}\fi
\noindent\textit{Source reference.} Original reference 08HD: sources/stacks-spaces-perfect/pinned-original-source.tex, source lines 3081--3081. Prerequisite pointer invoked at 3174; it names a standard fact of the same chapter that the counted argument uses, and every step of the counted proof built on it is retained verbatim in the delivered body.\par\medskip
\ifdefined\hypertarget\hypertarget{sourcepointer-3}{}\fi
\noindent\textit{Source reference.} Original reference 08H8: sources/stacks-spaces-perfect/pinned-original-source.tex, source lines 330--330. Prerequisite pointer invoked at 3161; it names a standard fact of the same chapter that the counted argument uses, and every step of the counted proof built on it is retained verbatim in the delivered body.\par\medskip
\ifdefined\hypertarget\hypertarget{sourcepointer-4}{}\fi
\noindent\textit{Source reference.} Original definition 08FT: sources/stacks-chow/tags.txt, source lines 10721--10721. Prerequisite pointer invoked at 3160; it names a standard fact of a sibling chapter that the counted argument uses, and every step of the counted proof built on it is retained verbatim in the delivered body.\par\medskip
\ifdefined\hypertarget\hypertarget{sourcepointer-5}{}\fi
\noindent\textit{Source reference.} Original lemma 08FV: sources/stacks-chow/tags.txt, source lines 10723--10723. Prerequisite pointer invoked at 3185; it names a standard fact of a sibling chapter that the counted argument uses, and every step of the counted proof built on it is retained verbatim in the delivered body.\par\medskip
\ifdefined\hypertarget\hypertarget{sourcepointer-6}{}\fi
\noindent\textit{Source reference.} Original lemma 08FX: sources/stacks-chow/tags.txt, source lines 10725--10725. Prerequisite pointer invoked at 3172; it names a standard fact of a sibling chapter that the counted argument uses, and every step of the counted proof built on it is retained verbatim in the delivered body.\par\medskip
\ifdefined\hypertarget\hypertarget{sourcepointer-7}{}\fi
\noindent\textit{Source reference.} Original lemma 08H4: sources/stacks-chow/tags.txt, source lines 10767--10767. Prerequisite pointer invoked at 3146; it names a standard fact of a sibling chapter that the counted argument uses, and every step of the counted proof built on it is retained verbatim in the delivered body.\par\medskip
\begin{lemma}
\label{lemma-descend-pseudo-coherent}
Let $X$ be a scheme. Let $E$ be an object of $D(\mathcal{O}_X)$.
The following are equivalent
\begin{enumerate}
\item $E$ is $m$-pseudo-coherent, and
\item $\epsilon^*E$ is $m$-pseudo-coherent on the small \'etale site of $X$.
\end{enumerate}
Here $\epsilon$ is as in (\ref{equation-epsilon}).
\end{lemma}

\begin{proof}
The implication (1) $\Rightarrow$ (2) is a general fact, see
Cohomology on Sites, Lemma
\ref{sites-cohomology-lemma-pseudo-coherent-pullback}.
Assume $\epsilon^*E$ is $m$-pseudo-coherent.
We will use without further mention that $\epsilon^*$ is
an exact functor and that therefore
$$
\epsilon^*H^i(E) = H^i(\epsilon^*E).
$$
To show that $E$ is $m$-pseudo-coherent we may work locally on $X$,
hence we may assume that $X$ is quasi-compact (for example affine).
Since $X$ is quasi-compact every \'etale covering $\{U_i \to X\}$
has a finite refinement. Thus we see that $\epsilon^*E$ is
an object of $D^{-}(\mathcal{O}_\etale)$, see
comments following
Cohomology on Sites, Definition
\ref{sites-cohomology-definition-pseudo-coherent}.
By Lemma \ref{lemma-epsilon-flat} it follows that $E$ is an object of
$D^-(\mathcal{O}_X)$.

\medskip\noindent
Let $n \in \mathbf{Z}$ be the largest integer such that
$H^n(E)$ is nonzero; then $n$ is also the largest integer
such that $H^n(\epsilon^*E)$ is nonzero.
We will prove the lemma by induction on $n - m$.
If $n < m$, then the lemma is clearly true.
If $n \geq m$, then $H^n(\epsilon^*E)$ is a finite
$\mathcal{O}_\etale$-module, see
Cohomology on Sites, Lemma \ref{sites-cohomology-lemma-finite-cohomology}.
Hence $H^n(E)$ is a finite $\mathcal{O}_X$-module, see
Lemma \ref{lemma-descend-finite-type}.
After replacing $X$ by the members of an open covering, we may
assume there exists a surjection $\mathcal{O}_X^{\oplus t} \to H^n(E)$.
We may locally on $X$ lift this to a map of complexes
$\alpha : \mathcal{O}_X^{\oplus t}[-n] \to E$ (details omitted).
Choose a distinguished triangle
$$
\mathcal{O}_X^{\oplus t}[-n] \to E \to C \to \mathcal{O}_X^{\oplus t}[-n + 1]
$$
Then $C$ has vanishing cohomology in degrees $\geq n$. On the other hand, the
complex $\epsilon^*C$ is $m$-pseudo-coherent, see
Cohomology on Sites, Lemma \ref{sites-cohomology-lemma-cone-pseudo-coherent}.
Hence by induction we see that $C$ is $m$-pseudo-coherent. Applying
Cohomology on Sites, Lemma \ref{sites-cohomology-lemma-cone-pseudo-coherent}
once more we conclude.
\end{proof}

\end{document}
